% !PRINTED 326
\weberpage{326}
Nehmen wir an, in (16) seien $n$, $n'$ selbst Diskriminanten:
\[
n = D_1, \quad n' = D_2, \quad D_1 \equiv D_2 \pmod{D},
\]
so ist
\[
(D, D_1) = (D, D_2)
\]
und es folgt aus (13):
\[
\begin{aligned}
(D, D_1) &= \pm(D_1, D) \\
(D, D_2) &= \pm(D_2, D),
\end{aligned}
\]
worin die oberen Zeichen gelten, wenn $D$ positiv ist, die unteren, wenn $D$, $D_1$, $D_2$ alle drei negativ sind, und verschiedene Zeichen, wenn $D$ negativ und $D_1$, $D_2$ von verschiedenen Vorzeichen sind. Es ist also nach (16):
\weq{17}{(D_1, D) = \pm(D_2, D),}
worin das untere Zeichen nur dann gilt, wenn die \gesperrt{beiden Zahlen} $D$ und $D_1 D_2$ beide negativ sind.

Sind $D_1$, $D_2$ ungerade, so können wir $D = 4m$ setzen, worin $m$ eine beliebige ganze Zahl ist, und wir erhalten aus (17)
\weq{18}{(D_1, m) = \pm(D_2, m),}
wenn
\weq{19}{D_1 \equiv D_2 \pmod{4m}}
\gesperrt{ist, und das untere Zeichen nur dann gilt, wenn $m$ und $D_1 D_2$ beide negativ sind.}

Dieser Satz ist zunächst bewiesen unter der Voraussetzung, daß $D_1$ und $D_2$ ungerade sind. Er gilt aber allgemein. Denn die Kongruenz (19) verlangt zunächst, daß $D_1$ und $D_2$ zugleich gerade oder zugleich ungerade seien. Sind sie beide gerade und ist dann $m$ gerade, so ist (18) richtig, weil beide Seiten verschwinden. Ist aber $m$ ungerade, $D_1$, $D_2$ gerade, so setze man in (17) $D = \pm m \equiv 1 \pmod{4}$, worin $m$ eine positive ungerade Zahl ist. Dann folgt aus $D_1 \equiv D_2 \pmod{D}$ die Kongruenz (19) und nach (17) ist
\[
(D_1, m) = (D_2, m),
\]
\weq{20}{\begin{aligned}
(D_1, m) &= (D_2, m),\\
(D_1, -m) &= \pm(D_2, -m), \\
(D_1, -1) &= \pm(D_2, -1),
\end{aligned}}
und folglich wieder
\[
(D_1, m) = (D_2, m).
\]
Also gilt diese Formel für jedes positive ungerade $m$, und nach (20) ist dann (18) auch für ein negatives $m$ erwiesen. Die Formeln (16) und (18) können zur Berechnung des Symbols $(D, n)$
%CONT>>
